\subsection{用装入记录快照检查一次正常传输}

下面选择四个时刻观察装入记录。设载荷长度为 \code{0x2100}，起点为
\code{0x00100000}。

\begin{longtable}{|L{0.19\linewidth}|L{0.21\linewidth}|L{0.22\linewidth}|L{0.26\linewidth}|}
\hline
观察时刻 & 下一写入地址 & 剩余字节 & 已经成立的事实 \\ \hline
头部通过后 & \code{0x00100000} & \code{0x2100} & 范围与入口合法，尚无载荷字节 \\ \hline
第 1 字节后 & \code{0x00100001} & \code{0x20ff} & 起始字节已写入并纳入校验 \\ \hline
第 64 字节后 & \code{0x00100040} & \code{0x20c0} & 第一块可向发送端确认 \\ \hline
最后字节后 & \code{0x00102100} & \code{0} & 长度满足，等待最终校验 \\ \hline
\end{longtable}

两个数值始终满足不变量：
\[
\texttt{destination\_next}-\texttt{load\_address}
=\texttt{image\_bytes}-\texttt{payload\_remaining}.
\]
左边是已经写入的地址跨度，右边是已经消费的字节数。若更新时漏掉递增或递减，其中一边会
立刻失配。固件可在调试版本中每轮检查该式。

\subsection{第 37 个字节损坏时哪些字节可见}

假设传输没有丢失长度，但第 37 个字节从 \code{0x11} 变成 \code{0x10}。装入器会继续收满
\code{0x2100} 字节，最终校验不匹配，状态转为 \code{REJECTED}。此时 RAM 中确实有一份
长度完整但内容错误的映像。

固件不必在发现校验失败时逐字节回滚，因为任务 PC 尚未提交。它必须保证两点：不把记录
转为 \code{READY}；下一次装入重新覆盖可能使用的范围。若机器要在多个已装入映像之间
保留旧版本，直接覆盖就不再足够，需要独立暂存区或双缓冲；本章只有一项任务，因此不引入
该结构。

\subsection{第 37 个字节之后串行溢出}

若控制器置 \code{RX\_OVERRUN}，固件知道至少一个字节丢失，却不知道具体有几个。此时
\code{destination\_next} 只表示已经消费的字节数，不能对应发送端位置。固件立即进入
\code{REJECTED}，丢弃直到下一次明确的传输同步标记，再请求整包重发。

仅把 \code{payload\_remaining} 继续减到零会造成错误对齐：后续字节向前填补缺口，头部
长度仍可能满足，但内容位置全部错位。校验大概率会失败，却要等整包结束才能发现。尽早
利用硬件溢出状态可以缩短失败路径。

\subsection{零填充区为什么不随载荷发送}

任务常需要一片初始值全为零的可写区域，例如计数器数组或暂存缓冲。如果发送端把几千个零
也逐字节发送，既增加传输时间，又没有增加信息。头部用 \code{zero\_bytes} 表达“映像末尾
之后这段范围必须被固件写零”。

这一做法建立了两类来源不同的任务字节：

\begin{longtable}{|L{0.24\linewidth}|L{0.26\linewidth}|L{0.36\linewidth}|}
\hline
区域 & 初值来源 & 校验与建立规则 \\ \hline
映像区 & 外部载荷 & 每个字节参与传输校验 \\ \hline
零填充区 & 固件写零 & 范围由已验证头部决定，不占载荷 \\ \hline
栈区 & 固件初始化 & 预置返回地址，其余按启动策略清零 \\ \hline
\end{longtable}

零填充必须发生在校验通过之后或之前均可，但在状态进入 \code{READY} 前必须完成。若清零
中途发生复位，启动流程重新开始，不能沿用旧的 \code{READY} 标记。

\subsection{为什么不能把整个 RAM 都清零}

全清零看似简单，却会覆盖固件工作区，包括当前装入记录和固件栈。即使先避开工作区，清零
整片 RAM 也会让启动时间随容量增长。按头部和固定布局只初始化任务真正依赖的区域，使
写入范围与所有权一致。

\subsection{任务栈如何支持一次真实函数调用}

求和核心本身是叶函数，但正式任务通常会调用子程序。设任务在写结果前调用位于
\code{0x00100100} 的格式化例程，调用指令地址为 \code{0x00100018}，顺序后继为
\code{0x0010001c}。调用前：

\begin{lstlisting}
PC = 0x00100018
SP = 0x001efffc
MEM32[0x001efffc] = 0x00000300   // outer return
\end{lstlisting}

\code{CALL} 先把 SP 降到 \code{0x001efff8}，写入内部返回地址
\code{0x0010001c}，再进入子程序：

\begin{lstlisting}
PC = 0x00100100
SP = 0x001efff8
MEM32[0x001efff8] = 0x0010001c   // return to task
MEM32[0x001efffc] = 0x00000300   // return to firmware
\end{lstlisting}

子程序的 \code{RET} 先消费 \code{0x0010001c}，恢复到任务内部，SP 回到
\code{0x001efffc}。任务最外层 \code{RET} 再消费 \code{0x00000300}，返回固件。两个
返回地址按后进先出次序工作。

\begin{center}
\begin{tikzpicture}[x=1cm,y=0.7cm]
\draw[BookRule,thick] (0,0) rectangle (8,4.4);
\draw[BookRule] (0,1.1) -- (8,1.1);
\draw[BookRule] (0,2.2) -- (8,2.2);
\draw[BookRule] (0,3.3) -- (8,3.3);
\node[anchor=west] at (0.2,3.85) {\code{0x001f0000}：栈区上界};
\node[anchor=west] at (0.2,2.75) {\code{0x001efffc}：返回固件 \code{0x00000300}};
\node[anchor=west] at (0.2,1.65) {\code{0x001efff8}：返回任务 \code{0x0010001c}};
\node[anchor=west] at (0.2,0.55) {更深调用与局部保存继续向低地址扩展};
\draw[flow] (9.3,1.65) -- (8.05,1.65);
\node[anchor=west] at (9.4,1.65) {子程序中的 SP};
\end{tikzpicture}
\end{center}
\diagramnote{返回地址的顺序来自 SP 的移动和 RAM 中的布局，不来自处理器对“函数层次”的
理解。破坏任一栈字都会改变后续控制流。}

\subsection{谁负责保证栈不越界}

当前处理器只检查地址是否位于整个 RAM，不知道任务栈下界
\code{0x001ec000}。若调用过深，SP 可以越过下界并覆盖其他任务数据，只要地址仍在 RAM，
硬件就会接受。固件给出 16 KiB 空间并不能强制任务停在边界内。

可以让编写任务的人证明最大栈深，或在每次函数入口加入软件检查；两者仍依赖任务合作。
真正的隔离需要地址访问权限，后续才会从任务相互覆盖的失败中推导。

\subsection{移交失败时由谁收束状态}

在 PC 提交之前，固件是唯一执行者，也是装入记录、任务目标区和串行协议的所有者。任何
失败都由固件把记录转为 \code{REJECTED}，发送错误码，并回到等待新头部状态。

PC 提交之后，固件不再执行。任务拥有当前寄存器与任务栈，直到 \code{RET}。此后出现的
无限循环、任意 RAM 写入和栈破坏，不能再由“装入器检查一次”解决。阶段不同，能采取行动
的主体也不同。

\begin{longtable}{|L{0.20\linewidth}|L{0.24\linewidth}|L{0.28\linewidth}|L{0.18\linewidth}|}
\hline
阶段 & 当前执行者 & 可修改的关键状态 & 失败收束 \\ \hline
接收头部 & 固件 & 装入记录、串行槽 & 拒绝传输 \\ \hline
接收载荷 & 固件 & 任务 RAM、校验状态 & 拒绝并重传 \\ \hline
建立现场 & 固件 & 任务栈、参数寄存器 & 不提交 PC \\ \hline
任务运行 & 任务 & 寄存器及全部可写物理 RAM & 合作返回或整机错误入口 \\ \hline
固定返回 & 固件 & 固件栈、发送状态 & 报告结果并等待 \\ \hline
\end{longtable}

这张所有权表给出了本章最重要的时间边界：同一片 RAM 在不同阶段由不同代码解释和修改，
但硬件尚未实施所有权。当前所有权是分析软件责任的工具，不是安全保证。
